%% kapa: exam=1 round=35 year=2024 period=2 form=A subject=appraisal_law subject_ko=감정평가관계법규 seq=15 number=15
%% source: Q-Net 감정평가사 기출문제 (cst00309) artlSeq=5213229 file=raw/1008/5213229/2교시 문제지 원본.pdf
%% conversion: kordoc 4.24.51 -> tools/gen1.py (auto); review=proofread-sonnet; flags=none
\begin{question}{15}
부동산 가격공시에 관한 법령상 비주거용 부동산가격의 공시에 관한 설명으로 옳지 않은 것은?
\begin{choices}
  \item 공시기준일 이후에 「건축법」에 따른 대수선이 된 비주거용 일반부동산은 해당 비주거용 개별부동산가격의 공시기준일을 다르게 할 수 있다.
  \item 비주거용 표준부동산의 임시사용승인일은 비주거용 표준부동산가격의 공시사항에 포함되지 않는다.
  \item 비주거용 표준부동산가격은 국가 등이 그 업무와 관련하여 비주거용 개별부동산가격을 산정하는 경우에 그 기준이 된다.
  \item 국토교통부장관은 비주거용 집합부동산가격을 공시하기 위하여 그 가격을 산정할 때에는 비주거용 집합부동산의 소유자와 그 밖의 이해관계인의 의견을 들어야 한다.
  \item 국토교통부장관은 공시한 비주거용 집합부동산가격의 오기를 정정하려는 경우에는 중앙부동산가격공시위원회의 심의를 거치지 아니할 수 있다.
\end{choices}
\end{question}
\answer{2}
